\documentclass[12pt]{article}
\usepackage[utf8]{inputenc}
\usepackage[T1]{fontenc}
\usepackage{hyperref}
\title{Old Latin Language Resources for K-12 Curriculum Development}
\author{Compiled by Claude Sonnet 5.5 (claude-sonnet-5-5) for Andr\'as Kornai}
\date{}
\begin{document}
\maketitle

\section{Introduction}
Old Latin (qbb) is the earliest attested stage of Latin, known only from inscriptions and fragments. It has no native speakers, so any curriculum would be a classical-studies or epigraphy course rather than a mother-tongue programme. The searches found no dedicated Old Latin corpus or school materials, only entries on individual inscriptions and an epigraphic photo collection.

\section{Primary Educational Resources}
No Old Latin textbooks or school programmes were found. A Textkit forum thread discusses the Praeneste fibula for learners.

\section{Linguistic and Reference Materials}
\begin{itemize}
\item Wikipedia, Duenos inscription (7th to 5th century BC; translation remains uncertain): \url{https://en.wikipedia.org/wiki/Duenos_inscription}
\item Wikipedia, Praeneste fibula (authenticity disputed): \url{https://en.wikipedia.org/wiki/Praeneste_fibula}
\item Britannica, Praeneste Fibula: \url{https://www.britannica.com/topic/Praeneste-Fibula}
\item Review of A. Gordon's 1975 monograph on the fibula's authenticity (PDF, in Portuguese, Humanitas 29--30): \url{https://ap1.sib.uc.pt/bitstream/10316.2/29320/1/humanitas29-30_artigo40.pdf}
\end{itemize}

\section{Religious and Community Resources}
None were identified.

\section{Digital and Audio Resources}
\begin{itemize}
\item Ohio State University Epigraphy collection, CIL I$^2$ 21--24 photographs (entries link to the Epigraphic Database Roma): \url{https://epigraphy.osu.edu/collections/latin-photo/cil-i2-21-24}
\item Textkit forum thread on the Praeneste fibula: \url{https://www.textkit.com/t/praeneste-fibula/15946}
\end{itemize}

\section{Curriculum Development Context}
Old Latin could be taught as a short unit within Latin or ancient-history courses, built around a few famous inscriptions and the Epigraphic Database Roma, with attention to open questions of dating, authenticity and translation.

\end{document}